
\section{Energy Storage Systems}

Energy Storage Systems (ESS) are the pivotal component that enables the decoupling of energy generation from consumption, thereby facilitating the effective utilization of intermittent renewable resources. Among the various storage technologies available, Lithium-Ion (Li-ion) batteries have emerged as the dominant technology for residential applications due to their high energy density, decreasing cost, and long cycle life \cite{Chen2009}.

However, Li-ion batteries are complex electrochemical devices subject to degradation over time. This degradation manifests as capacity fade (loss of energy storage capability) and power fade (increase in internal resistance), both of which ultimately reduce the economic viability of the storage system. Degradation is typically categorized into two modes. Calendar aging occurs invariably over time, regardless of usage, and is primarily influenced by the state of charge (SoC) and temperature. Cycling aging, on the other hand, is driven by the charge and discharge current throughput and the depth of discharge (DoD) \cite{Xu2018}.

Optimizing the operation of a BESS requires a careful balance. While aggressive cycling might capture more short-term arbitrage revenue, it accelerates degradation, potentially incurring replacement costs that outweigh the operational gains. Therefore, sophisticated EMS strategies must explicitly account for battery health. Recent studies have integrated degradation cost models into the optimization objective, demonstrating that health-aware control strategies significantly extend battery life and improve the overall Net Present Value (NPV) of the project \cite{Grimaldi2024}. This thesis adopts a simplified yet effective degradation penalty term in the reward function to internalize this long-term cost into the real-time decision-making process.
